% ---------------------------------------------------------------------------
%  Praeambel
% ---------------------------------------------------------------------------
% Zeichenkodierung und Schrift engine-abhaengig:
%   pdfLaTeX  -> inputenc/fontenc + Latin Modern (8 Bit)
%   XeLaTeX / LuaLaTeX -> UTF-8 direkt ueber fontspec.  Wichtig: unter XeLaTeX
%   darf fontenc T1 NICHT geladen werden, sonst wird "ss" falsch gesetzt.
\usepackage{iftex}
\ifPDFTeX
  \usepackage[utf8]{inputenc}
  \usepackage[T1]{fontenc}
  \usepackage{lmodern}
\else
  \usepackage{fontspec}
\fi
\usepackage[ngerman]{babel}
\usepackage[german=quotes]{csquotes}
\usepackage{microtype}

\usepackage{amsmath,amssymb}
\usepackage{graphicx}
\usepackage{booktabs}
\usepackage{longtable}
\usepackage{tabularx}
\usepackage{multirow}
\usepackage{array}
\usepackage{xcolor}
\usepackage{caption}
\usepackage{subcaption}
\usepackage{enumitem}
\usepackage{float}
\usepackage{setspace}
\usepackage{tikz}
\usetikzlibrary{arrows.meta,positioning,fit,backgrounds,calc,shapes.geometric,
                patterns,decorations.pathreplacing}

\usepackage[a4paper,inner=3.2cm,outer=2.6cm,top=2.8cm,bottom=3.0cm]{geometry}


\usepackage[T1]{fontenc}
\usepackage[utf8]{inputenc}
%\usepackage[german]{babel}
%\usepackage{lmodern}
%\usepackage{microtype}
%\usepackage{amsmath,amssymb,amsthm}
%\usepackage{booktabs}
%\usepackage{tabularx}
%\usepackage{enumitem}
%\usepackage{csquotes}
%\usepackage[hidelinks]{hyperref}

\numberwithin{equation}{chapter}
\setcounter{chapter}{2}   % naechstes \chapter wird Kapitel 3

% Notationskonventionen
\newcommand{\Dom}[1]{\mathcal{D}_{#1}}          % Verteilung
\newcommand{\Hcls}{\mathcal{H}}                  % Hypothesenklasse
\newcommand{\HdH}{\Hcls\Delta\Hcls}              % symmetrische Differenzklasse
\newcommand{\IoU}{\mathrm{IoU}}
\DeclareMathOperator*{\argmin}{arg\,min}

% ---------------------------------------------------------------------------
%  Literatur.  backend=biber ist der Standardweg und wird von latexmk
%  automatisch bedient (biber liest die .bcf-Datei).
%  Ohne biber im System: backend=bibtex setzen UND latexmk mit
%    -e '$bibtex_use=2;' aufrufen, sonst startet latexmk trotzdem biber.
% ---------------------------------------------------------------------------
\usepackage[backend=biber,style=authoryear,maxcitenames=2,mincitenames=1,
            maxbibnames=12,giveninits=true,dashed=false,sorting=nyt,
            urldate=long]{biblatex}
\addbibresource{literatur.bib}
\DeclareNameAlias{sortname}{family-given}
\renewcommand*{\nameyeardelim}{\addcomma\space}   % (Autor, 2016)

\usepackage[hidelinks]{hyperref}
\hypersetup{
  pdftitle={Datenaugmentierung fuer die CNN-basierte Detektion von Rohrprofil-Stirnflaechen},
  pdfsubject={Bachelorarbeit},
  colorlinks=true, linkcolor=black, citecolor=black, urlcolor=black,
  linktoc=all
}

% ---------------------------------------------------------------------------
%  Satzspiegel und Typografie
% ---------------------------------------------------------------------------
\onehalfspacing
\setlength{\emergencystretch}{3em}
\captionsetup{font=small,labelfont=bf,width=0.94\textwidth}
\captionsetup[sub]{font=small}
\setlist{itemsep=2pt,parsep=0pt,topsep=4pt}

% Kopfzeilen
\usepackage[automark,headsepline]{scrlayer-scrpage}
\pagestyle{scrheadings}
\clearpairofpagestyles
\ihead*{\headmark}
\ohead*{\pagemark}

% Trennhilfen fuer die wiederkehrenden langen Komposita
\hyphenation{Trans-la-ti-ons-äqui-va-ri-anz Trans-la-ti-ons-in-va-ri-anz
  Augmen-tie-rungs-pipe-line Augmen-tie-rung Ob-jekt-aus-schnit-te
  Merk-mals-ex-trak-ti-on Sicht-bar-keits-rech-nung Ge-samt-schritt-wei-te
  De-tek-ti-ons-kopf De-tek-ti-ons-ebe-ne In-fe-renz-auf-lö-sung
  Ob-jekt-grö-ßen-ver-tei-lung Trai-nings-ver-tei-lung Kon-fi-denz-in-ter-vall
  Rohr-pro-fil-Stirn-flä-chen Anti-Ab-kür-zungs-Maß-nah-me
  Ein-bet-tungs-rea-lis-mus Prä-dik-ti-ons-po-si-tio-nen}

% ---------------------------------------------------------------------------
%  Eigene Makros
% ---------------------------------------------------------------------------
\newcommand{\code}[1]{\texttt{\small #1}}          % Dateien, Parameter, Pfade
\newcommand{\px}{\,px}
\newcommand{\PP}{\,Prozentpunkte}
\newcommand{\mAPfifty}{mAP@0{,}5}
\newcommand{\mAPseven}{mAP@0{,}75}
\newcommand{\mAPall}{mAP@0{,}5{:}0{,}95}
\newcommand{\Oe}{\ensuremath{\varnothing}}         % Durchmesser als Laengeneinheit

% Absatzueberschrift im Stil des Theorieteils
\newcommand{\absatz}[1]{\textbf{#1}}

% Tabellenhilfen
\newcolumntype{L}[1]{>{\raggedright\arraybackslash}p{#1}}
\newcolumntype{R}[1]{>{\raggedleft\arraybackslash}p{#1}}
\newcommand{\tabkopf}[1]{\textbf{#1}}
